\section*{Background and Related Work}
\label{sec:background}

\subsection*{UAV Path Planning}
\label{subsec:uav_path_planning}

Unmanned Aerial Vehicles (UAVs) have become essential for autonomous navigation in complex environments with static and dynamic obstacles. Path planning for UAVs involves finding collision-free trajectories from a start position to a goal while satisfying constraints such as terrain clearance, obstacle avoidance, and energy efficiency \cite{Aggarwal2023, Zhang2023}.

Traditional path planning algorithms such as Rapidly-exploring Random Trees (RRT) and A* provide fundamental solutions but often struggle with dynamic environments and real-time constraints. RRT-based approaches explore configuration space probabilistically and can handle high-dimensional spaces efficiently \cite{Lavalle1998}, while A* offers guaranteed optimality in discretized environments but requires complete prior knowledge of the environment \cite{Hart1968}.

However, these classical methods have limitations: they lack adaptive behavior in changing environments, often produce suboptimal paths, and require significant computational resources for online replanning. Consequently, recent research has shifted toward using Particle Swarm Optimization (PSO) for UAV path planning because PSO naturally balances exploration and exploitation while being computationally efficient \cite{Bonyadi2015, Kennedy1995}.

Furthermore, deep reinforcement learning approaches have shown promise for UAV navigation, with agents learning adaptive policies directly from environment observations \cite{Schaal2018, Lillicrap2015}. These methods enable UAVs to generalize across different environments and handle dynamic obstacles more effectively than hand-crafted policies.

\subsection*{Particle Swarm Optimization}
\label{subsec:pso_fundamentals}

PSO is a swarm-based metaheuristic algorithm where each particle in the swarm represents a solution to the optimization problem. PSO updates particle velocities and positions at each time step, based on three components: (1) the best position a specific particle has found (cognitive component), (2) the best position found by the entire swarm (social component), and (3) the previous velocity, governed by the inertia weight (exploration-exploitation balance) \cite{Kennedy1995, Bonyadi2015}.

Formally, the velocity and position update equations of PSO are:

\begin{equation}
\begin{cases}
v_i^{(t+1)} &= w v_i^{(t)} + c_1 r_1^{(t)} \odot (p_i^{(t)} - x_i^{(t)}) + c_2 r_2^{(t)} \odot (g^{(t)} - x_i^{(t)}) \\
x_i^{(t+1)} &= x_i^{(t)} + v_i^{(t+1)}
\end{cases}
\label{eq:pso_update}
\end{equation}

where:
\begin{itemize}
\item $w$ is the inertia weight controlling the influence of previous velocity (exploration vs. exploitation trade-off),
\item $c_1$ and $c_2$ are acceleration coefficients governing cognitive and social components, respectively,
\item $r_1^{(t)}, r_2^{(t)} \sim \mathcal{U}(0,1)^D$ are random vectors providing stochastic exploration,
\item $v_i^{(t)}$ and $x_i^{(t)}$ are velocity and position of particle $i$ at time step $t$,
\item $p_i^{(t)}$ is the personal best position found by particle $i$,
\item $g^{(t)}$ is the global best position found by the entire swarm,
\item $\odot$ denotes element-wise multiplication.
\end{itemize}

\subsubsection*{PSO Parameter Tuning}

Early research by Shi and Eberhart \cite{ShiEberhart1998} validated that $w = 0.729844$, $c_1 = c_2 = 1.496180$ satisfy theoretical convergence conditions and provide good balance for most problems. However, these fixed parameters are suboptimal for many real-world problems. Research has shown that adjusting these parameters to manage exploration and exploitation improves PSO convergence and solution quality \cite{Bonyadi2015}.

Traditional schedules for parameter adaptation include:

\begin{itemize}
\item \textbf{Progress Definition}: $p(t) = \frac{t}{T}$, where $t \in \{1, \ldots, T\}$ and $T$ is the maximum number of iterations.

\item \textbf{Inertia Weight Schedule (Linearly Decreasing)}:
\begin{equation}
w(t) = w_{\max} - (w_{\max} - w_{\min}) \cdot p(t)
\end{equation}
with typical defaults: $w_{\max} = 0.9$, $w_{\min} = 0.1$ \cite{ShiEberhart1998}.

\item \textbf{Time-Varying Acceleration Coefficients (TVAC)}:
\begin{align}
c_1(t) &= c_{1,\max} - (c_{1,\max} - c_{1,\min}) \cdot p(t) \\
c_2(t) &= c_{2,\min} + (c_{2,\max} - c_{2,\min}) \cdot p(t)
\end{align}
with typical bounds: $c_{1,\min} = 0.5$, $c_{1,\max} = 2.5$, $c_{2,\min} = 0.5$, $c_{2,\max} = 2.5$ \cite{Ratnaweera2004}.
\end{itemize}

These deterministic schedules enforce early exploration (high $w, c_1$, low $c_2$) and gradual shift to exploitation (low $w, c_1$, high $c_2$). However, such fixed schedules are problem-dependent and may not adapt well to the dynamics of the optimization landscape.

\subsection*{Parameter Adaptation Strategies in PSO}
\label{subsec:parameter_adaptation}

The standard PSO suffers from limitations including premature convergence (convergence to local optima due to particles clustering around a suboptimal region) and sensitivity to initial parameters \cite{Bonyadi2015}. To address these issues, four main improvement strategies have emerged: (1) modification of PSO control parameters, (2) hybridization with other metaheuristics (e.g., genetic algorithms, differential evolution), (3) cooperation and multi-swarm techniques, and (4) learning-based adaptive parameter control \cite{Bonyadi2015}.

Self-Adaptive Particle Swarm Optimization (SAPSO) methods attempt to adjust control parameters during optimization without introducing additional hyperparameters \cite{Engelbrecht2010}. Early approaches used fuzzy logic and memory-based heuristics to modify parameters online. However, these methods typically require explicit rules or assumptions about the optimization landscape.

More recently, reinforcement learning has emerged as a promising approach for online parameter adaptation. Unlike traditional schedules or heuristic rules, RL-based methods can learn adaptive policies directly from the optimization process \cite{Zhang2023RL}. These approaches treat parameter adaptation as a sequential decision problem: at each iteration, an RL agent observes the current state of the swarm (e.g., diversity, convergence progress) and selects appropriate parameter values to maximize overall performance \cite{Engelbrecht2023}.

The foundational work by Von Eschwege and Engelbrecht \cite{VonEschwege2024} introduced SAC-SAPSO (Soft Actor-Critic Self-Adaptive PSO), which combines Soft Actor-Critic deep reinforcement learning with PSO parameter adaptation. This approach demonstrated that RL-based parameter adaptation can achieve 50-80\% improvement in solution quality over baseline PSO while being robust and time-invariant with minimal tunable hyperparameters.

\subsection*{Reinforcement Learning for Optimization}
\label{subsec:reinforcement_learning}

Reinforcement learning (RL) offers a principled framework for learning adaptive control policies through interaction with an environment. In the context of optimization, the environment is the target optimization problem (e.g., PSO dynamics), and the RL agent learns to control algorithm parameters to maximize a cumulative reward signal (e.g., fitness improvement) \cite{Haarnoja2018, Lillicrap2015}.

\subsubsection*{Actor-Critic Methods}

Actor-Critic algorithms, a major class of policy-gradient methods, maintain two components: an actor (policy network) that selects actions and a critic (value network) that evaluates state quality \cite{Konda2000, Mnih2016}. This architecture provides stable learning with reduced variance compared to pure policy-gradient methods \cite{Schulman2016}.

The original DDPG (Deep Deterministic Policy Gradient) algorithm \cite{Lillicrap2015} extended actor-critic methods to continuous control by using deterministic policies and experience replay, enabling learning in high-dimensional action spaces. However, DDPG suffered from instability due to overestimation of Q-values and difficult hyperparameter tuning \cite{Fujimoto2018}.

\subsubsection*{Soft Actor-Critic (SAC)}

Soft Actor-Critic \cite{Haarnoja2018, Haarnoja2019} addresses DDPG's limitations through entropy regularization, automatic temperature tuning, and a maximum-entropy RL objective. The SAC framework combines:

\begin{itemize}
\item \textbf{Off-policy Learning}: Uses experience replay buffer to sample past experiences, improving sample efficiency.
\item \textbf{Twin Q-Networks}: Maintains two independent Q-function approximators to reduce overestimation bias.
\item \textbf{Entropy Regularization}: Encourages exploration by adding entropy bonus $\alpha H(\pi(\cdot|s))$ to the reward, where $H$ is policy entropy and $\alpha$ is automatically tuned.
\item \textbf{Stochastic Policy}: Learns a probabilistic policy $\pi$ instead of deterministic actions, naturally balancing exploration and exploitation.
\end{itemize}

The SAC objective maximizes:
\begin{equation}
J(\pi) = \mathbb{E}_{t \sim D}[R(s_t, a_t) + \alpha H(\pi(\cdot|s_t))]
\end{equation}
where $D$ is the replay buffer, $R$ is the cumulative reward, and $\alpha$ is automatically adjusted via dual optimization \cite{Haarnoja2018}.

SAC has demonstrated superior performance on continuous control benchmarks (robotics, autonomous navigation) with better sample efficiency and stability than DDPG and PPO \cite{Haarnoja2019}.

\subsubsection*{CrossQ Optimization}

Recent advances in off-policy learning include CrossQ \cite{Bhatt2019}, which improves sample efficiency by:

\begin{itemize}
\item \textbf{Batch Normalization}: Applies BatchNorm layers in critic networks to handle distribution mismatch between current and next state batches.
\item \textbf{No Target Networks}: Removes target networks (a major source of computational overhead) by leveraging BatchNorm's distribution properties.
\item \textbf{Unit Temporal Difference (UTD) Ratio}: Achieves state-of-the-art sample efficiency with just 1 critic update per environment sample (vs. 20 in prior work).
\end{itemize}

CrossQ maintains the core benefits of SAC while reducing computational cost, making it suitable for parameter adaptation in optimization algorithms where evaluation budgets are constrained.

\subsubsection*{State Representation and Attention Mechanisms}

For complex optimization problems, effective state representations are crucial for RL agents to understand the problem landscape and make informed decisions. Recent work leverages \textbf{attention mechanisms} to create problem-aware state encodings \cite{Vaswani2017}.

The Transformer architecture \cite{Vaswani2017} introduced self-attention, enabling networks to weight different parts of the input based on relevance. In the context of PSO parameter adaptation, attention can capture temporal dynamics of swarm behavior and identify critical features (e.g., convergence stagnation, diversity loss) that should influence parameter updates \cite{Vaswani2017}.

Multi-head attention with temporal windows enables the RL agent to:
\begin{itemize}
\item Capture long-range dependencies in swarm trajectories,
\item Identify patterns of premature convergence or excessive exploration,
\item Dynamically weight swarm features (diversity, fitness improvement rate, stagnation) based on optimization stage.
\end{itemize}

\subsubsection*{Reward Design}

The retained controller uses a direct reward tied only to global-best fitness improvement. This keeps the learning signal aligned with the optimization target and avoids additional reward-shaping terms that require separate tuning.

\begin{equation}
R(s, a) =
\begin{cases}
1, & \text{if global-best fitness improves} \\
-1, & \text{otherwise}
\end{cases}
\end{equation}

where each component addresses a specific optimization goal. This enables the RL agent to develop nuanced policies that avoid local optima, maintain healthy swarms, and achieve fast convergence \cite{Pitis2019}.

\subsection*{Summary and Motivation}
\label{subsec:summary_motivation}

Traditional PSO with fixed or heuristically-scheduled parameters limits performance across diverse problem instances. While SAPSO methods have improved upon standard PSO, they typically rely on hand-crafted rules or simple adaptation mechanisms. The recent SAC-SAPSO work \cite{VonEschwege2024} demonstrated promising results by combining deep RL with PSO, achieving significant performance improvements.

However, existing approaches have not fully explored:

\begin{itemize}
\item \textbf{Fine-grained parameter granularity}: Most work adapts global parameters or simple subgroups. Per-particle adaptation enables more nuanced control suited to heterogeneous swarm behaviors.
\item \textbf{Advanced state representations}: While basic metrics (diversity, fitness) are useful, attention-based temporal modeling of swarm dynamics can capture richer behavioral patterns.
\item \textbf{Sample-efficient learning}: Optimization problems typically have limited evaluation budgets. Combining state-of-the-art RL (SAC with CrossQ) with efficient state encodings can improve learning with fewer PSO evaluations.
\item \textbf{Unified framework}: Integration of direct fitness-improvement rewards, attention mechanisms, transformer-based state encoding, and advanced RL algorithms into a cohesive system remains underexplored.
\end{itemize}

Our AFSACPSO (Advanced Parameter Exploration CrossQ-SAC for PSO) algorithm addresses these gaps by:

\begin{enumerate}
\item \textbf{Per-Particle Parameter Control}: Learns individual inertia and acceleration coefficients for each particle, enabling adaptive heterogeneous behavior within the swarm.
\item \textbf{Transformer-Based State Encoding}: Uses multi-head self-attention over a temporal window of swarm metrics to create rich, context-aware state representations.
\item \textbf{SAC with CrossQ Optimization}: Leverages the latest advances in deep RL for stable, sample-efficient learning with automatic entropy tuning.
\item \textbf{Multi-Objective Reward Design}: Balances fitness improvement, diversity maintenance, convergence speed, and curiosity through weighted reward composition.
\item \textbf{BatchNorm-Enhanced Networks}: Applies batch normalization in critic networks as per CrossQ, improving stability and reducing computational overhead.
\end{enumerate}

This integrated approach enables AFSACPSO to automatically discover problem-specific parameter adaptation strategies that outperform fixed schedules and heuristic rules across diverse optimization benchmarks.

% References
\begin{thebibliography}{99}

\bibitem{Aggarwal2023}
Aggarwal, M., Kumar, A. (2023).
``Autonomous localized path planning algorithm for UAVs based on TD3 strategy,''
\textit{Scientific Reports}, 14, 51349.

\bibitem{Zhang2023}
Zhang, Y., Liu, J., et al. (2023).
``Autonomous Navigation of the UAV through Deep Reinforcement Learning with Sensor Perception Enhancement,''
\textit{Mathematical Problems in Engineering}.

\bibitem{Lavalle1998}
LaValle, S.~M. (1998).
``Rapidly-exploring random trees: A new tool for path planning,''
\textit{Technical Report}, Iowa State University.

\bibitem{Hart1968}
Hart, P.~E., Nilsson, N.~J., Raphael, B. (1968).
``A formal basis for the heuristic determination of minimum cost paths,''
\textit{IEEE Transactions on Systems Science and Cybernetics}, 4(2), 100--107.

\bibitem{Kennedy1995}
Kennedy, J., Eberhart, R.~C. (1995).
``Particle swarm optimization,''
\textit{Proceedings of the IEEE International Conference on Neural Networks}, pp. 1942--1948.

\bibitem{Bonyadi2015}
Bonyadi, M.~R., Michalewicz, Z. (2015).
``Particle Swarm Optimization: A survey of historical and recent developments with hybridization perspectives,''
\textit{arXiv:1804.05319}.

\bibitem{Schaal2018}
Schaal, S., Peters, J., Nakanishi, J. (2018).
``Learning robot control with recurrent neural networks,''
\textit{International Journal of Robotics Research}.

\bibitem{Lillicrap2015}
Lillicrap, T.~P., Hunt, J.~J., Pritzel, A., Heess, N., Erez, T., Tassa, Y., Silver, D., Wierstra, D. (2015).
``Continuous control with deep reinforcement learning,''
\textit{arXiv:1509.02971}.

\bibitem{Konda2000}
Konda, V.~R., Tsitsiklis, J.~N. (2000).
``Actor-critic algorithms,''
\textit{SIAM Journal on Control and Optimization}, 42(4), 1143--1166.

\bibitem{Mnih2016}
Mnih, V., Badia, A.~P., Mirza, M., Graves, A., Lillicrap, T., Harley, T., Silver, D., Kavukcuoglu, K. (2016).
``Asynchronous methods for deep reinforcement learning,''
\textit{ICML}, pp. 1928--1937.

\bibitem{Schulman2016}
Schulman, J., Moritz, P., Levine, S., Jordan, M., Abbeel, P. (2016).
``High-dimensional continuous control using generalized advantage estimation,''
\textit{ICLR}.

\bibitem{Fujimoto2018}
Fujimoto, S., van Hoof, H., Meger, D. (2018).
``Addressing function approximation error in actor-critic methods,''
\textit{ICML}, pp. 1587--1596.

\bibitem{Haarnoja2018}
Haarnoja, T., Zhou, A., Abbeel, P., Levine, S. (2018).
``Soft Actor-Critic: Off-Policy Maximum Entropy Deep Reinforcement Learning with a Stochastic Actor,''
\textit{ICML}, pp. 1861--1870.

\bibitem{Haarnoja2019}
Haarnoja, T., Zhou, A., Finn, C., Abbeel, P., Levine, S. (2019).
``Soft Actor-Critic Algorithms and Applications,''
\textit{ICML}.

\bibitem{Bhatt2019}
Bhatt, A., Palenicek, D., Belousov, B., Argus, M., Amiranashvili, A., Brox, T., Peters, J. (2019).
``CrossQ: Batch Normalization in Deep Reinforcement Learning for Greater Sample Efficiency and Simplicity,''
\textit{ICLR}, pp. 1909--1922.

\bibitem{Vaswani2017}
Vaswani, A., Shazeer, N., Parmar, N., Uszkoreit, J., Jones, L., Gomez, A.~N., Kaiser, Ł., Polosukhin, I. (2017).
``Attention Is All You Need,''
\textit{NeurIPS}, pp. 5998--6008.

\bibitem{Pitis2019}
Pitis, R. (2019).
``Generalizing Across Multi-Objective Reward Functions in Deep Reinforcement Learning,''
\textit{arXiv:1809.06364}.

\bibitem{VonEschwege2024}
Von Eschwege, D.~H., Engelbrecht, A.~P. (2024).
``Soft Actor-Critic Approach to Self-Adaptive Particle Swarm Optimisation,''
\textit{Mathematics}, 12(22), 3481.

\bibitem{Engelbrecht2010}
Engelbrecht, A.~P. (2010).
``Heterogeneous particle swarm optimization,''
\textit{Proceedings of the International Conference on Swarm Intelligence}.

\bibitem{Engelbrecht2023}
Engelbrecht, A.~P. (2023).
``Self-adaptive learning in particle swarm optimization,''
\textit{Machine Learning: Algorithms and Applications}.

\bibitem{ShiEberhart1998}
Shi, Y., Eberhart, R.~C. (1998).
``A modified particle swarm optimizer,''
\textit{Proceedings of the IEEE International Conference on Evolutionary Computation}, pp. 69--73.

\bibitem{Ratnaweera2004}
Ratnaweera, A., Halgamuge, S.~K., Watson, H.~C. (2004).
``Self-organizing hierarchical particle swarm optimizer with time-varying acceleration coefficients,''
\textit{IEEE Transactions on Evolutionary Computation}, 8(3), 240--255.

\bibitem{Zhang2023RL}
Zhang, Q., Zhou, A., Zhao, S., et al. (2023).
``Reinforcement-learning-based parameter adaptation method for particle swarm optimization,''
\textit{Complex \& Intelligent Systems}, 9, 1012--1025.

\end{thebibliography}

